% demo-full-report.tex -- a complete report using every command of the class
\documentclass[
  text={fontsize=11pt, lang=british},
  paths={code=code/},
  position={figure=H, table=H, program=H, pseudocode=H},
  sectioning={part=break, chapter=nobreak, annex=nobreak, bib=nobreak},
  toc={depth=1, annexdepth=0, bibdepth=0},
]{myreport}
\usepackage{csquotes}
\usepackage{clrscode3e}
\usepackage[locale=UK]{siunitx}

\mycourse*{Fluid mechanics}[MECA0025]
\myaddprofessor{Professor}{Paul}
\myaddassistant*{Assistant}{Alice}
\mytitle{Demonstration report}
\mysubtitle{Every command of myreport in one document}
\mylogo{square.jpg}
\myaddauthor{Author}{First}[s000001]
\myaddauthor*{Author}{Second}[s000002]

\begin{document}
\mymaketitle[2]

\myabstract[abstract]
This document uses every command of the class once, so that a change in
any of them is visible at a glance.

\mytableofcontents
\mylistoffigures
\mylistoftables
\mylistofprograms
\mylistofequations

\myintroduction[number=1, label=intro]
The introduction is numbered as the first part; see \myref{part:theory}.

\mypart{Theory}[label=theory]
\mychapter{Definitions}[label=defs]
\mysection{Dimensionless numbers}[label=numbers]
\begin{myequation}[re]{Reynolds number}
  Re = \frac{\rho U L}{\mu} = \frac{\num{1.2} \cdot \num{14.4} \cdot \num{2.41}}{\num{1.9e-5}}
\end{myequation}
\mysubsection{Courant number}[label=cfl]
\begin{myequation}[cfl]{Courant number}
  CFL &= \frac{U \Delta t}{\Delta x} \\
  \Delta t &= \frac{CFL \cdot \Delta x}{U}
\end{myequation}
\mysubsubsection{Remark}
See \myseeref{eq:re} and \myref{subsec:cfl}.

\mypart{Simulations}[label=sims]
\mychapter{Set-up}[label=setup]
\myfigure{wide.png}[Domain]{The simulation domain. A second sentence that
does not go to the list of figures.}
\begin{mysub}{mesh}{The mesh and a zoom on it.}
  \mysubfigure{wide.png}[Mesh]{0.55}\hfill
  \mysubfigure{square.jpg}[Zoom]{0.35}
\end{mysub}
\begin{mytable}[firstrow=2]{|l|c|}{params}{Simulation parameters.}
  \hline \textbf{Parameter} & \textbf{Value} \\ \hline
  Velocity & \qty{14.4}{\meter\per\second} \\
  CFL & \num{4} \\
  Duration & \num{8} FTT \\ \hline
\end{mytable}
\mychapter{Code}[label=code]
\myprogram{hello}{hello.py}{python}{A post-processing script.}
\begin{mypseudocode}{search}{Searching for the optimal angle.}
\proc{Optimise}($\theta_{min}, \theta_{max}$)
\li \For $\theta \gets \theta_{min}$ \To $\theta_{max}$
\li   \Do run a simulation
    \End
\li \Return the best angle
\end{mypseudocode}
Results: \myref{fig:wide}, \myref{fig:mesh:b}, \myref{tab:params},
\myref{code:py:hello}, \myref{code:pseudo:search}, cited from
\cite{art1,book1,web1}.

\mypart{Conclusion}
The conclusion. See also \myref{ann:part:details}.

\begin{myannex}
  \mypart{Details}[label=details]
  \mychapter{Long derivation}
  Text of the annex.
\end{myannex}

\mybibliography
\end{document}
